\documentclass[../../../include/open-logic-section]{subfiles}

\begin{document}

\olfileid{sth}{choice}{countablechoice}
\olsection{الاختيار المعدود}

من اليسير، من غير استعانة بالاختيار ولا بالترتيب الجيد، إثبات الآتي.

\begin{lem}[في $\Zminus$]
لكل مجموعة منتهية دالة اختيار.
\end{lem}

\begin{proof}
لتكن~${\OLId{latin-ordinary}{a}} = \{{\OLId{latin-ordinary}{b}}_{\OLMathNumeral{1}}, \ldots, {\OLId{latin-ordinary}{b}}_{\OLId{latin-ordinary}{n}}\}$. ولنفترض، تيسيرًا، أن كل
${\OLId{latin-ordinary}{b}}_{\OLId{latin-ordinary}{i}} \neq \emptyset$. فثمة أشياء~${\OLId{latin-ordinary}{c}}_{\OLMathNumeral{1}}, \ldots, {\OLId{latin-ordinary}{c}}_{\OLId{latin-ordinary}{n}}$ بحيث
${\OLId{latin-ordinary}{c}}_{\OLMathNumeral{1}} \in {\OLId{latin-ordinary}{b}}_{\OLMathNumeral{1}}, \ldots, {\OLId{latin-ordinary}{c}}_{\OLId{latin-ordinary}{n}} \in {\OLId{latin-ordinary}{b}}_{\OLId{latin-ordinary}{n}}$. ثم تضمن
\olref[z][pairs]{prop:pairsconsequences} وجود المجموعة
$\{\langle {\OLId{latin-ordinary}{b}}_{\OLMathNumeral{1}}, {\OLId{latin-ordinary}{c}}_{\OLMathNumeral{1}}\rangle , \ldots, \langle {\OLId{latin-ordinary}{b}}_{\OLId{latin-ordinary}{n}}, {\OLId{latin-ordinary}{c}}_{\OLId{latin-ordinary}{n}}\rangle\}$؛
وهي دالة اختيار لـ~${\OLId{latin-ordinary}{a}}$.
\end{proof}

فإذا انتقلنا إلى المجموعات غير المنتهية لم يبق الأمر على هذا الظهور.
وتأمل أضعف امتداد طبيعي لما سبق.

\begin{defish}
\emph{الاختيار المعدود.} لكل مجموعة \emph{معدودة} دالة اختيار.
\end{defish}

وهذه صورة خاصة من الاختيار. وقد تبين أن الرياضيين كثيرًا ما استعملوها
من غير أن ينبهوا على موضع استعمالها. وإليك مثالين حسنين.
\footnote{نقلهما~\citet[\S9.4]{Potter2004} ولوكا إنكورفاتي.}

\begin{ex}
من المعاني القريبة إلى الذهن أنه لكل مجموعة~${\OLId{latin-looped}{A}}$، إما
$\cardle{\omega}{{\OLId{latin-looped}{A}}}$، أو~$\cardeq{{\OLId{latin-looped}{A}}}{{\OLId{latin-ordinary}{n}}}$ لبعض~${\OLId{latin-ordinary}{n}} \in \omega$. وهذا
تقرير للقول إن كل مجموعة إما منتهية وإما غير منتهية. وقد ذكر كانتور،
وغيره من الرياضيين، هذا الحكم من غير برهان. وأما نحن فأثبتناه في
\olref[cardinals][classing]{generalinfinitycharacter}، غير أن عملنا
هنالك كان داخل~$\ZFC$: افترضنا إمكان ترتيب كل مجموعة~${\OLId{latin-looped}{A}}$ ترتيبًا
جيدًا، فضُمِن وجود~$\card{{\OLId{latin-looped}{A}}}$. فكان الاختيار مفروضًا صراحة.

وقد أقام~\citet{Dedekind1888} برهانه الخاص على الحكم على هذا الوجه.

\begin{thm}[في $\Zminus + \text{الاختيار المعدود}$]
لكل~${\OLId{latin-looped}{A}}$، إما~$\cardle{\omega}{{\OLId{latin-looped}{A}}}$ أو~$\cardeq{{\OLId{latin-looped}{A}}}{{\OLId{latin-ordinary}{n}}}$ لبعض
${\OLId{latin-ordinary}{n}} \in \omega$.
\end{thm}

\begin{proof}
افترض~$\cardneq{{\OLId{latin-looped}{A}}}{{\OLId{latin-ordinary}{n}}}$ لكل~${\OLId{latin-ordinary}{n}} \in \omega$. وحينئذ، لكل
${\OLId{latin-ordinary}{n}} < \omega$، توجد مجموعة جزئية~${\OLId{latin-looped}{A}}_{\OLId{latin-ordinary}{n}} \subseteq {\OLId{latin-looped}{A}}$ عدد عناصرها بالضبط
$\cardexpo{{\OLMathNumeral{2}}}{{\OLId{latin-ordinary}{n}}}$. وبالمتتالية~${\OLId{latin-looped}{A}}_{\OLMathNumeral{0}}, {\OLId{latin-looped}{A}}_{\OLMathNumeral{1}}, {\OLId{latin-looped}{A}}_{\OLMathNumeral{2}}, \ldots$ نعرف لكل~${\OLId{latin-ordinary}{n}}$:
\[
	{\OLId{latin-looped}{B}}_{\OLId{latin-ordinary}{n}} = {\OLId{latin-looped}{A}}_{\OLId{latin-ordinary}{n}} \setminus \bigcup_{{\OLId{latin-ordinary}{i}} < {\OLId{latin-ordinary}{n}}} {\OLId{latin-looped}{A}}_{\OLId{latin-ordinary}{i}}.
\]
ثم لاحظ:
\begin{align*}
	\card{\bigcup_{{\OLId{latin-ordinary}{i}} < {\OLId{latin-ordinary}{n}}}{\OLId{latin-looped}{A}}_{\OLId{latin-ordinary}{i}}}
	&\leq \card{{\OLId{latin-looped}{A}}_{\OLMathNumeral{0}}} + \card{{\OLId{latin-looped}{A}}_{\OLMathNumeral{1}}} + \ldots + \card{{\OLId{latin-looped}{A}}_{{\OLId{latin-ordinary}{n}}-{\OLMathNumeral{1}}}}\\
	&={\OLMathNumeral{1}} + {\OLMathNumeral{2}} + \ldots + {\OLMathNumeral{2}}^{{\OLId{latin-ordinary}{n}}-{\OLMathNumeral{1}}}\\
	& = {\OLMathNumeral{2}}^{\OLId{latin-ordinary}{n}} - {\OLMathNumeral{1}}\\
	& < {\OLMathNumeral{2}}^{\OLId{latin-ordinary}{n}} = \card{{\OLId{latin-looped}{A}}_{\OLId{latin-ordinary}{n}}}
\end{align*}
فلكل~${\OLId{latin-looped}{B}}_{\OLId{latin-ordinary}{n}}$ عنصر على الأقل~${\OLId{latin-ordinary}{c}}_{\OLId{latin-ordinary}{n}}$. وهذه المجموعات~${\OLId{latin-looped}{B}}_{\OLId{latin-ordinary}{n}}$ منفصلة اثنتين
اثنتين؛ فإذا~${\OLId{latin-ordinary}{c}}_{\OLId{latin-ordinary}{n}} = {\OLId{latin-ordinary}{c}}_{\OLId{latin-ordinary}{m}}$، كان~${\OLId{latin-ordinary}{n}} = {\OLId{latin-ordinary}{m}}$. ولما كان كل~${\OLId{latin-ordinary}{c}}_{\OLId{latin-ordinary}{n}} \in {\OLId{latin-looped}{A}}$،
كانت الدالة~${\OLId{latin-ordinary}{f}}({\OLId{latin-ordinary}{n}}) = {\OLId{latin-ordinary}{c}}_{\OLId{latin-ordinary}{n}}$ حقنًا~$\omega \to {\OLId{latin-looped}{A}}$.
\end{proof}
\noindent
لم يصرح ديديكند بأنه استعمل الاختيار المعدود. أفبان لك موضعه؟ فأعد
النظر في البرهان، ثم أعده مرة أخرى.

استعمل الاختيار المعدود مرتين: أولاهما لتحصيل متتالية المجموعات~${\OLId{latin-looped}{A}}_{\OLMathNumeral{0}}$
و${\OLId{latin-looped}{A}}_{\OLMathNumeral{1}}$ و${\OLId{latin-looped}{A}}_{\OLMathNumeral{2}}$ و\dots\@، والأخرى لاختيار~${\OLId{latin-ordinary}{c}}_{\OLId{latin-ordinary}{n}}$ من كل~${\OLId{latin-looped}{B}}_{\OLId{latin-ordinary}{n}}$. وليس
هذا الاستعمال مما يحذف؛ فقد أثبت~\citet[p.~138]{Cohen1966} أن الحكم
يفشل إذا عدمنا كل صورة من صور الاختيار. أي يتسق مع~$\ZF$ أن توجد
مجموعات لا تقارن بـ~$\omega$.
\end{ex}

\begin{ex}
في~\citeyear{Cantor1878} قرر كانتور أن الاتحاد المعدود للمجموعات
المعدودة معدود، ولم يذكر برهانًا، ولعله استغنى بظهوره. وقد أثبتنا صورة
أعم في~\olref[card-arithmetic][simp]{kappaunionkappasize}، لكن برهانها
صرح بفرض الاختيار. بل حتى الصورة الأخص تحتاج إلى الاختيار المعدود.

\begin{thm}[في $\Zminus + \text{الاختيار المعدود}$]
إذا كانت~${\OLId{latin-looped}{A}}_{\OLId{latin-ordinary}{n}}$ معدودة لكل~${\OLId{latin-ordinary}{n}} \in \omega$، فإن
$\bigcup_{{\OLId{latin-ordinary}{n}} < \omega} {\OLId{latin-looped}{A}}_{\OLId{latin-ordinary}{n}}$ معدودة.
\end{thm}

\begin{proof}
افترض، بلا إخلال بالعموم، أن كل~${\OLId{latin-looped}{A}}_{\OLId{latin-ordinary}{n}} \neq \emptyset$. فعندئذ توجد
لكل~${\OLId{latin-ordinary}{n}} \in \omega$ !!a{surjection}
${\OLId{latin-ordinary}{f}}_{\OLId{latin-ordinary}{n}} \colon \omega \to {\OLId{latin-looped}{A}}_{\OLId{latin-ordinary}{n}}$. عرّف
${\OLId{latin-ordinary}{f}} \colon \omega \times \omega \to \bigcup_{{\OLId{latin-ordinary}{n}} < \omega} {\OLId{latin-looped}{A}}_{\OLId{latin-ordinary}{n}}$
بـ~${\OLId{latin-ordinary}{f}}({\OLId{latin-ordinary}{m}}, {\OLId{latin-ordinary}{n}}) = {\OLId{latin-ordinary}{f}}_{\OLId{latin-ordinary}{n}}({\OLId{latin-ordinary}{m}})$. والنتيجة لازمة لأن
$\omega \times \omega$ معدودة
(\olref[sfr][siz][zigzag]{natsquaredenumerable}) ولأن~${\OLId{latin-ordinary}{f}}$
!!a{surjection}.
\end{proof}
\noindent
وموضع الاختيار المعدود هنا تعيين متتالية الدوال~${\OLId{latin-ordinary}{f}}_{\OLMathNumeral{0}}$ و${\OLId{latin-ordinary}{f}}_{\OLMathNumeral{1}}$ و${\OLId{latin-ordinary}{f}}_{\OLMathNumeral{2}}$
و\dots\footnote{وقع استخدام مشابه
للاختيار في~\olref[card-arithmetic][simp]{kappaunionkappasize} عندما
أصدرنا التعليمة «لكل~$\beta \in \cardfont{{\OLId{latin-ordinary}{a}}}$، ثبّت
!!a{injection}~${\OLId{latin-ordinary}{f}}_\beta$».} ولا تستغني النتيجة عنه هي أيضًا. فقد أثبت
\citet{FefermanLevy1963} أنه يتسق مع~$\ZF$ أن تكون كاردينالية اتحاد
معدود من مجموعات معدودة هي~$\beth_{\OLMathNumeral{1}}$. وقد عبر راسل عن المعنى بصورة
أطرف:
\begin{quote}
  يوضح هذا مليونير كان يشتري زوجًا من الجوارب كلما اشترى زوجًا من
  الأحذية، ولا يشتريه في أي وقت آخر، وكان مولعًا بشراء كليهما إلى حد
  أنه امتلك في النهاية~$\aleph_{\OLMathNumeral{0}}$ زوجًا من الأحذية و$\aleph_{\OLMathNumeral{0}}$ زوجًا
  من الجوارب\dots\@ ويمكننا في الأحذية التمييز بين الأيمن والأيسر،
  ولذلك نستطيع أن نختار واحدًا من كل زوج، أي نستطيع اختيار الأحذية
  اليمنى كلها أو اليسرى كلها؛ أما الجوارب فلا يوحي لنا مبدأ اختيار
  مماثل بنفسه، ولا نستطيع أن نضمن، ما لم نفترض المسلّمة الضربية [أي
  الاختيار في الواقع]، وجود أي فئة تتألف من جورب واحد من كل زوج.
  \citep[p.~126]{Russell1919}
\end{quote}
فخلاصة المعنى أن ضربًا من الاختيار لازم لإثبات أن من ملك عددًا معدودًا
من أزواج الجوارب لم يملك إلا عددًا معدودًا من الجوارب. ومن غير
الاختيار المعدود، أو ما يكافئه، قد لا يكون الاتحاد المعدود لمجموعات
معدودة معدودًا.
\end{ex}

والحاصل أن الاختيار المعدود دخل في أعمال كثيرة من غير تنبه ظاهر من
مستعمليه. والسؤال الفلسفي هو: بم نسوغ هذا الأصل؟

وقد يلتمس له تسويغ تصوري بمهمة فائقة: نفرض أن الاختيار الأول يقع في
$\nicefrac{{\OLMathNumeral{1}}}{{\OLMathNumeral{2}}}$ دقيقة، والثاني في~$\nicefrac{{\OLMathNumeral{1}}}{{\OLMathNumeral{4}}}$ دقيقة،
و\dots، والاختيار رقم~${\OLId{latin-ordinary}{n}}$ في~$\nicefrac{{\OLMathNumeral{1}}}{{\OLMathNumeral{2}}^{\OLId{latin-ordinary}{n}}}$ دقيقة، و\dots\@
فنكون قد أنجزنا، في~${\OLMathNumeral{1}}$ دقيقة، متتالية من الاختيارات نمطها~$\omega$
وحددنا دالة اختيار.

لكن ماذا تثبت تجربة فكرية كهذه؟ إنها أولًا تحمل معنى «الاختيار» على
حقيقته الحسية حملًا شديدًا؛ ثم إنها تربط صدق الرياضيات بالإمكان
الميتافيزيقي.

وفوق ذلك كله، لا تفيد تسويغ الاختيار مطلقًا، بل الاختيار المعدود وحده.
فإذا أردنا دالة اختيار لكل مجموعة، لزم أن نستطيع إتمام «مهمة فائقة ذات
طول ترتيبي كيفما كان». وهذا، عند التحقيق، مما يدعو إلى الهزء.

\end{document}
